\section{Quantitative results}
\label{sec:results}

\TODO{update the status line when the evaluation is complete} The retrieval comparison and the
threshold sweep make no model call and are complete on the original 30 questions. The evaluation
with the language models (\code{evaluate.py run --judge}) is spread over several days by the
50-request daily quota and was still running when this version was written. The tables below are
generated from the result files by the report's script; a red TODO stands where a table has not been
generated yet.

\subsection{Corpus and question set}

\begin{table}[htbp]
  \centering\small
  \begin{tabularx}{\textwidth}{@{}X l@{}}
    \toprule
    Item & Value \\
    \midrule
    Documents & 50: 20 WHO guidelines (PDF), 30 PubMed Central articles (text) \\
    Records read by \code{load_documents} (2026-11-14) & 1{,}530, of which 1{,}500 WHO pages \\
    Chunks in the index & 3,410 \\
    Evidence quotes found on their page (\code{check_eval_set.py}) & 20 / 20 \\
    Questions & 40: A01--A20 answerable, O01--O10 out of corpus, N01--N10 held out \\
    \bottomrule
  \end{tabularx}
  \caption{Corpus and question set.}
  \label{tab:corpus}
\end{table}

Each answerable question has a ground-truth answer, a source document, a page and a verbatim
evidence quote. The quotes come from 19 of the 50 documents, 8 of them PubMed articles. The
out-of-corpus questions are seven medical questions and three general-knowledge ones; the held-out
questions are described in Section~\ref{sec:difficulties}.

\subsection{Retrieval and gates, with no model call}

On the answerable questions, hybrid retrieval finds the source document for all 20 (hit rate 1.0)
and the evidence page for 18 (page hit rate 0.90), against 0.95 and 0.80 for dense retrieval
alone, with a context precision of 0.911 against 0.857. BM25 alone finds the most pages (0.95) but
cannot gate. Figure~\ref{fig:separation} shows, for each of the 30 questions, the two scores the
gates use. Both scores put every answerable question above every out-of-corpus one; the dense
similarity leaves a gap of 0.036 around its threshold (0.6235 against 0.5878), the cross-encoder a
gap from 0.156 to 0.870 around 0.5. The \code{heldout_auroc} and \code{heldout_margin} columns need a
new run of \code{evaluate.py retrieval} on the 40 questions. \TODO{rerun \code{evaluate.py retrieval}
and report the held-out separation}

\IfFileExists{\gen/retrieval_table.tex}{\input{\gen/retrieval_table.tex}}{\TODO{retrieval table: \code{docs/report/generated/retrieval_table.tex}, from \code{results/retrieval.csv}}}

\begin{figure}[htbp]
  \centering
  \IfFileExists{\res/retrieval_questions.csv}{%
  \begin{tikzpicture}
    \begin{axis}[width=0.72\textwidth, height=0.36\textwidth,
                 xlabel={best dense similarity (BGE-M3)},
                 ylabel={MiniLM gate score (sigmoid)},
                 xmin=0.3, xmax=0.85, ymin=-0.03, ymax=1.03,
                 legend pos=south east, legend style={font=\scriptsize},
                 tick label style={font=\scriptsize}, label style={font=\small}]
      \addplot[only marks, mark=*, mark size=1.8pt, blue!70!black]
        table[col sep=comma, x=dense_best,
              y expr={1/(1+exp(-\thisrow{hybrid+minilm_gate_best}))},
              skip coords between index={20}{1000}] {\res/retrieval_questions.csv};
      \addlegendentry{answerable (A01--A20)}
      \addplot[only marks, mark=x, mark size=2.6pt, thick, red!75!black]
        table[col sep=comma, x=dense_best,
              y expr={1/(1+exp(-\thisrow{hybrid+minilm_gate_best}))},
              skip coords between index={0}{20}, skip coords between index={30}{1000}] {\res/retrieval_questions.csv};
      \addlegendentry{out of corpus (O01--O10)}
      \draw[dashed, gray] (axis cs:0.6,-0.03) -- (axis cs:0.6,1.03);
      \draw[dashed, gray] (axis cs:0.3,0.5) -- (axis cs:0.85,0.5);
      \node[font=\scriptsize, gray, anchor=south west] at (axis cs:0.6,0.02) {0.6};
      \node[font=\scriptsize, gray, anchor=south west] at (axis cs:0.31,0.5) {0.5};
    \end{axis}
  \end{tikzpicture}}{\TODO{plot of \code{results/retrieval_questions.csv}: file missing}}
  \caption{Best score of each of the original 30 questions under the two gates, read from
  \code{results/retrieval_questions.csv}. The cross-encoder's logit is mapped to 0..1 with a sigmoid,
  as \code{rerank.gate_score} does. Dashed lines: the two thresholds. The isotretinoin question (O06)
  is the cross nearest to 0.6.}
  \label{fig:separation}
\end{figure}

\subsection{Threshold sweep}

\code{evaluate.py sweep} scores 11 thresholds of the dense similarity, from 0.30 to 0.80, on the
original 30 questions (table generated from \code{results/threshold_sweep.csv}). At 0.60 the answer rate and the refusal rate are both
1.0; at 0.55 only 7 of the 10 out-of-corpus questions are refused, and at 0.65 one answerable
question is refused. So 0.6, the brief's value, is also the best on these questions; since it was
checked on the same 30 questions, this score is optimistic.

\IfFileExists{\gen/sweep_table.tex}{\input{\gen/sweep_table.tex}}{\TODO{sweep table: \code{docs/report/generated/sweep_table.tex}, from \code{results/threshold_sweep.csv}}}



\subsection{Latency}

Table~\ref{tab:latency} gives the measured times. Retrieval and the gate take well under a second;
the model call dominates the time of an answer.

\begin{table}[htbp]
  \centering\small
  \begin{tabularx}{\textwidth}{@{}X r@{}}
    \toprule
    Step (CPU, per question unless stated) & Time \\
    \midrule
    Dense retrieval, top 5 & 185\,ms \\
    BM25, top 5 & 7\,ms \\
    Hybrid retrieval (dense 20 + BM25 20, fused) & 0.22\,s \\
    Hybrid + MiniLM gate, five passages in one batch (comparison) & 0.8\,s \\
    Hybrid + MiniLM gate with early stop (app): answerable / out of corpus & about 0.3\,s / 0.5\,s \\
    Gate with MedCPT / bge-reranker-v2-m3 & about 3\,s / 11\,s \\
    Reordering 20 candidates: MiniLM / MedCPT / bge-reranker-v2-m3 & 2.2\,s / 14.7\,s / about 46\,s \\
    Model call (mean latency per answer, \code{summary.csv}) & \TODO{from summary table} \\
    Embedding the corpus with BGE-M3 (once) & about 3\,s per chunk, about 3\,h \\
    \bottomrule
  \end{tabularx}
  \caption{Measured times. Retrieval rows from \code{results/retrieval.csv} and the README; the
  early-stop row from commit \code{174cef4}.}
  \label{tab:latency}
\end{table}

\subsection{Answers with the language models}

Five systems run with each model: \code{rag_dense} (dense retrieval, similarity gate),
\code{rag_hybrid} (with BM25 fused in), \code{rag_rerank} (hybrid passages, cross-encoder gate; the
app's default), and the two baselines the course requires, \code{bm25_rag} (5 BM25 passages, same
prompt and citation check, no gate) and \code{no_rag} (the model alone). The metrics are the answer
rate on A01--A20, the refusal rate on O01--O10, the refusal rate on N01--N10, hit rate, context
precision, faithfulness judged by the default model, the number of judged answers and the mean
latency. The plan in \code{scripts/run_evaluation.py} is the five systems with
\nolinkurl{nvidia/nemotron-3-ultra-550b-a55b:free}, then \code{rag_rerank} with
\nolinkurl{google/gemma-4-31b-it:free}.

\IfFileExists{\gen/summary_table.tex}{\input{\gen/summary_table.tex}}{\TODO{summary table: \code{docs/report/generated/summary_table.tex}, from \code{results/summary.csv}}}

\TODO{two or three sentences on the summary table once the run is complete: refusal of the held-out
questions per gate, faithfulness, \code{no_rag} against the RAG systems, the second model}

\TODO{human ratings: mean \code{correct_0_1}, \code{safety_1_5} and \code{clarity_1_5} from
\nolinkurl{results/human_rating_sheet.csv}, with the number of rated rows}

\begin{figure}[htbp]
  \centering
  \fbox{\parbox[c][2.2cm][c]{0.8\textwidth}{\centering\TODO{screenshot of the app answering A06, with
  the sources table (\code{07_EXTRA/screenshots/})}}}
  \caption{The Gradio app.}
  \label{fig:app}
\end{figure}
